\begin{proposition}[局部联络的规范变换]
设主丛局部截面换为 $s^\prime=sg$，联络形式的局部代表为 $A=s^*\boldsymbol\omega$。则 $A^\prime=\operatorname{Ad}_{g^{-1}}A+g^{-1}\dd g$；矩阵群情形第一项为 $g^{-1}Ag$。
\end{proposition}
\begin{proof}
换局部截面 $s'=sg$ 时，新截面沿 $X$ 的切向量 $s'_*X$ 既含旧截面切向量的右平移，又含 $g$ 自身变化产生的垂直分量。联络形式对水平部分用 $G$-等变性、对垂直部分用再现性质，两项合起来给出非齐次变换律。

令 $s'=sg$，其中 $g:U\to G$。对 $X\in T_xM$，先计算新截面沿 $X$ 的切向量 $s'_*X$。由乘积截面微分，$s'_*X$ 同时含有 $s_*X$ 经右平移到 $s(x)g(x)$ 的水平部分，以及 $g(x)$ 自身沿 $X$ 变化产生的垂直部分。把后者经右平移回到单位元再经 基本 竖直 场 嵌入，得
\begin{equation*}
s'_*X
=(R_g)_*s_*X+\bigl(g^{-1}\dd g(X)\bigr)_P(sg),
\end{equation*}
其中 $(R_g)_*$ 是右平移 $R_g:h\mapsto hg$ 的切映射，$\bigl(g^{-1}\dd g(X)\bigr)_P$ 是把 $\mathfrak g$ 元素 $g^{-1}\dd g(X)$ 经 基本 场 映射 $\mathfrak g\to T_{sg}P$ 嵌入所得垂直向量。

由 $A'(X)=\boldsymbol\omega_{s'(x)}(s'_*X)$ 且 $s'(x)=s(x)g(x)$，代入上式并利用联络形式的线性性，
\begin{equation*}
A'(X)
=\boldsymbol\omega_{sg}\bigl((R_g)_*s_*X\bigr)
 +\boldsymbol\omega_{sg}\bigl(\bigl(g^{-1}\dd g(X)\bigr)_P(sg)\bigr).
\end{equation*}

对第一项用 $G$-等变性 $\boldsymbol\omega_{sg}\bigl((R_g)_*v\bigr)=\operatorname{Ad}_{g^{-1}}\boldsymbol\omega_s(v)$，取 $v=s_*X$，
\begin{equation*}
\boldsymbol\omega_{sg}\bigl((R_g)_*s_*X\bigr)
=\operatorname{Ad}_{g^{-1}}\boldsymbol\omega_s(s_*X)
=\operatorname{Ad}_{g^{-1}}A(X).
\end{equation*}
$G$-等变性是联络形式的定义公理之一，几何意义是右作用把联络形式按伴随表示拉回。

对第二项用再现性质 $\boldsymbol\omega_p\bigl(\xi_P(p)\bigr)=\xi$，取 $\xi=g^{-1}\dd g(X)$、$p=sg$，
\begin{equation*}
\boldsymbol\omega_{sg}\bigl(\bigl(g^{-1}\dd g(X)\bigr)_P(sg)\bigr)
=g^{-1}\dd g(X).
\end{equation*}
再现性质是联络形式的另一条定义公理：在基本竖直场上联络形式是恒等映射。

合并两项。
\begin{equation*}
A'(X)
=\operatorname{Ad}_{g^{-1}}A(X)+g^{-1}\dd g(X).
\end{equation*}
对矩阵 Lie 群 $\operatorname{Ad}_{g^{-1}}A=g^{-1}Ag$，于是
\begin{equation*}
A'(X)=g^{-1}A(X)g+g^{-1}\dd g(X),
\end{equation*}
即式\eqnrefs{eq:principal-local-connection-transform}。非齐次项来自局部截面本身也在变化，所以把联络系数当作普通张量会漏掉这一项。

典型例子。
\begin{itemize}
\item $G=U(1)$：$\operatorname{Ad}_{g^{-1}}=\mathrm{id}$（Abel 群 伴随 平凡），$A'=A+\ii\,\dd\chi$（取 $g=\ee^{\ii\chi}$）。这正是电磁势的规范变换 $A_\mu\mapsto A_\mu+\partial_\mu\chi$。
\item $G=SU(N)$：$A'=g^{-1}Ag+g^{-1}\dd g$，非 Abel 规范变换。$g^{-1}Ag$ 是共轭转动，$g^{-1}\dd g$ 是纯规范补偿。在 Yang--Mills 理论中，这一变换律保证 $D\psi=\dd\psi+A\psi$ 与 $\psi$ 同步齐次变换 $\psi\mapsto g^{-1}\psi$。
\end{itemize}

非齐次项 $g^{-1}\dd g$ 是规范势在规范变换下的"补偿"，对应物理上规范势 $A_\mu$ 在波函数相位重定向下获得的导数项。它使规范势不能视为张量，但保证协变导数 $D\psi$ 是张量（齐次变换）。这一非齐次性是规范理论全部相互作用结构的根源。
\end{proof}

\begin{proposition}[联络的换标架律]
对向量丛局部标架 $e^\prime=e\Lambda$，若 $\nabla e_b=e_a\otimes\omega^a{}_b$，则 $\omega^\prime=\Lambda^{-1}\omega\Lambda+\Lambda^{-1}\dd\Lambda$。
\end{proposition}
\begin{proof}
联络由 $\nabla e_c=e_a\otimes\omega^a{}_c$ 定义。换标架 $e'=e\Lambda$ 后，对 $\nabla e'_b$ 用 Leibniz 展开再换回新基，读出新系数即得非齐次变换律。

设新标架 $e'_b=e_c\Lambda^c{}_b$，其中 $\Lambda:U\to GL(n,\RR)$ 光滑。由联络对被求导向量满足 Leibniz 规则，对 $\nabla e'_b=\nabla(e_c\Lambda^c{}_b)$ 逐项展开，
\begin{align*}
\nabla e'_b
&=\nabla(e_c\Lambda^c{}_b)\\
&=(\nabla e_c)\Lambda^c{}_b+e_c\otimes\dd\Lambda^c{}_b.
\end{align*}
第一项是"基向量被联络作用"再乘系数，第二项是"系数被外微分"乘基向量——这正是联络作为方向导数的 Leibniz 性质。

由旧标架下联络一形式的定义 $\nabla e_c=e_a\otimes\omega^a{}_c$，代回第一项，
\begin{equation*}
(\nabla e_c)\Lambda^c{}_b
=e_a\otimes\omega^a{}_c\Lambda^c{}_b.
\end{equation*}

把指标 $c$ 在第一项中作哑指标，$a$ 作自由指标，合并两项
\begin{align*}
\nabla e'_b
&=e_a\otimes\omega^a{}_c\Lambda^c{}_b
 +e_c\otimes\dd\Lambda^c{}_b\\
&=e_a\otimes\bigl(\omega^a{}_c\Lambda^c{}_b+\dd\Lambda^a{}_b\bigr),
\end{align*}
其中第二行把第二项的哑指标 $c$ 改写为 $a$。

为读出新标架下的联络一形式 $\omega'$，需把基向量 $e_a$ 换回 $e'_r$。由 $e'_b=e_c\Lambda^c{}_b$ 取逆得 $e_a=e'_r(\Lambda^{-1})^r{}_a$，代回
\begin{align*}
\nabla e'_b
&=e'_r(\Lambda^{-1})^r{}_a\otimes
 \bigl(\omega^a{}_c\Lambda^c{}_b+\dd\Lambda^a{}_b\bigr)\\
&=e'_r\otimes
 \Bigl[
  (\Lambda^{-1})^r{}_a\omega^a{}_c\Lambda^c{}_b
  +(\Lambda^{-1})^r{}_a\dd\Lambda^a{}_b
 \Bigr].
\end{align*}

与新标架下定义 $\nabla e'_b=e'_r\otimes\omega'^r{}_b$ 比较基向量 $e'_r$ 的系数，得
\begin{equation*}
\omega'^r{}_b
=(\Lambda^{-1})^r{}_a\omega^a{}_c\Lambda^c{}_b
+(\Lambda^{-1})^r{}_a\dd\Lambda^a{}_b,
\end{equation*}
即\eqnrefs{eq:math-dg-connection-gauge-transform}。第一项是齐次的共轭转动，第二项是非齐次的纯规范补偿。

典型例子。
\begin{itemize}
\item 正交换标架 $\Lambda\in SO(n)$：$\Lambda^{-1}=\Lambda^T$，$\omega'=\Lambda^T\omega\Lambda+\Lambda^T\dd\Lambda$。第二项取值于 $\mathfrak{so}(n)$（反对称），保持度量相容性。
\item 局部 Lorentz 标架（四脚标架）变换 $\Lambda\in SO(1,3)$：$\omega'=\Lambda^{-1}\omega\Lambda+\Lambda^{-1}\dd\Lambda$，正是广义相对论中 自旋联络 在局部 Lorentz 变换下的规律。
\item 复标架变换 $\Lambda\in U(N)$：$\omega'=\Lambda^\dagger\omega\Lambda+\Lambda^\dagger\dd\Lambda$，对应 $U(N)$ 规范势的规范变换。
\end{itemize}

非齐次项 $\Lambda^{-1}\dd\Lambda$ 是"联络不是张量"的精确含义。在广义相对论中，它保证协变导数 $\nabla_\mu V^\nu=\partial_\mu V^\nu+\Gamma^\nu{}_{\mu\rho}V^\rho$ 在坐标变换下张量性成立——$\Gamma$ 的非齐次变换正好抵消 $\partial_\mu V^\nu$ 中多出的二阶项。在规范理论中，它使规范势 $A_\mu$ 在规范变换下获得 $\partial_\mu\chi$ 补偿项，从而协变导数 $D_\mu\psi$ 齐次变换。
\end{proof}

\begin{lemma}[协变外微分的平方]
设 $D$ 为向量丛联络诱导的协变外微分，曲率为 $\Omega$。对任意丛值微分形式 $\eta$，有 $D^2\eta=\Omega\wedge\eta$，右边使用曲率的纤维上的线性作用。
\end{lemma}
\begin{proof}
由协变外微分的定义 $D\eta=\dd\eta+\omega\wedge\eta$，其中 $\omega$ 是取值于 $\End E$ 的联络一形式、$\eta$ 是 $E$-值 $p$-形式。再次作用 $D$，先把 $D\eta$ 视作新的 $E$-值 $(p+1)$-形式代入同一公式，
\begin{equation*}
D^2\eta
=D(D\eta)
=\dd(D\eta)+\omega\wedge(D\eta).
\end{equation*}
把 $D\eta=\dd\eta+\omega\wedge\eta$ 代入两处，
\begin{equation*}
D^2\eta
=\dd(\dd\eta+\omega\wedge\eta)
 +\omega\wedge(\dd\eta+\omega\wedge\eta).
\end{equation*}
由外微分的线性性拆开，
\begin{equation*}
D^2\eta
=\dd^2\eta+\dd(\omega\wedge\eta)
 +\omega\wedge\dd\eta+\omega\wedge(\omega\wedge\eta).
\end{equation*}
由 $\dd^2=0$ 第一项消失。对第二项用分次 Leibniz 规则 $\dd(\alpha\wedge\beta)=\dd\alpha\wedge\beta+(-1)^{|\alpha|}\alpha\wedge\dd\beta$，其中 $|\omega|=1$，
\begin{equation*}
\dd(\omega\wedge\eta)
=\dd\omega\wedge\eta+(-1)^1\omega\wedge\dd\eta
=\dd\omega\wedge\eta-\omega\wedge\dd\eta.
\end{equation*}
代回，
\begin{equation*}
D^2\eta
=\dd\omega\wedge\eta-\omega\wedge\dd\eta
 +\omega\wedge\dd\eta+\omega\wedge\omega\wedge\eta.
\end{equation*}
中间两项 $-\omega\wedge\dd\eta+\omega\wedge\dd\eta$ 严格相消，只剩
\begin{equation*}
D^2\eta
=\dd\omega\wedge\eta+\omega\wedge\omega\wedge\eta
=(\dd\omega+\omega\wedge\omega)\wedge\eta.
\end{equation*}
由曲率二形式定义 $\Omega:=\dd\omega+\omega\wedge\omega$，代回得
\begin{equation*}
D^2\eta=\Omega\wedge\eta,
\end{equation*}
即\eqnrefs{eq:math-dg-covariant-derivative-square}。这也解释了为何曲率必然以 $\dd\omega+\omega\wedge\omega$ 而不是单独 $\dd\omega$ 出现。
\end{proof}

\begin{proposition}[曲率的有限扰动]
若 $\widetilde\nabla=\nabla+A$ 且 $A\in\Omega^1(M;\End E)$，则 $\widetilde R=R+D_\nabla A+A\wedge A$。此式是有限差的精确恒等式，未截断扰动阶。
\end{proposition}
\begin{proof}
曲率是联络的二次函数（含 $\nabla^2$ 和 $\nabla\wedge\nabla$），故联络的有限扰动 $A$ 产生线性项（$D_\nabla A$）和二次项（$A\wedge A$），结构完全由交换子代数决定。

把新联络写成
\begin{equation*}
 \widetilde\nabla_X=\nabla_X+A_X,
 \qquad
 A_X:=A(X,\cdot).
\end{equation*}
$A$ 是 $\Omega^1(M;\operatorname{End}E)$ 中任一元素，$\widetilde\nabla$ 与 $\nabla$ 同为 $E$ 上联络当且仅当它们的差是张量一形式。

曲率是交换子减联络沿括号方向的作用：
\begin{equation*}
 \widetilde R(X,Y)
 =[\widetilde\nabla_X,\widetilde\nabla_Y]-\widetilde\nabla_{[X,Y]}
 =[\nabla_X+A_X,\nabla_Y+A_Y]
 -(\nabla_{[X,Y]}+A_{[X,Y]}).
\end{equation*}

按双线性展开：
\begin{itemize}
\item 纯 $\nabla$ 部分：$[\nabla_X,\nabla_Y]-\nabla_{[X,Y]}=R(X,Y)$，即旧曲率。
\item 一次 $A$ 的交叉项：$\nabla_XA_Y-\nabla_YA_X-A_{[X,Y]}$，组合成协变外微分 $(D_\nabla A)(X,Y)$，即 $A$ 沿旧联络的协变外微分。
\item 二次项：$[A_X,A_Y]=A_XA_Y-A_YA_X$，即 $A\wedge A$。
\end{itemize}
逐项核对：$[\nabla_X,A_Y]=\nabla_XA_Y-A_Y\nabla_X$ 中 $-A_Y\nabla_X$ 与 $[A_X,\nabla_Y]=A_X\nabla_Y-\nabla_YA_X$ 中的 $A_X\nabla_Y$ 不直接相消，但与 $\nabla_XA_Y-\nabla_YA_X$ 合并后正好是协变外微分 $D_\nabla A$ 的定义。

因此
\begin{equation*}
 \widetilde R=R+D_\nabla A+A\wedge A,
\end{equation*}
与\eqnrefs{eq:math-dg-curvature-finite-perturbation}完全相同。曲率对联络"线性化 + 自相互作用"的结构因此不是坐标巧合，而是联络仿射几何的必然结果：线性项 $D_\nabla A$ 是一阶变分（用于 Chern--Weil 理论和指标定理），二次项 $A\wedge A$ 是非 Abel 自相互作用（与 Yang--Mills 方程的非线性性同源）。

取 $A_t=tA$，$\widetilde R_t=R+tD_\nabla A+t^2 A\wedge A$。对 $\operatorname{tr}(\widetilde R_t^k)$ 在 $t=0$ 求导，只有线性项 $D_\nabla A$ 贡献（二次项 $t^2$ 在 $t=0$ 导数为零）。这正是一阶变分控制 Chern--Weil 超渡 的原因：特征类的变分只看 $D_\nabla A$，非 Abel 自相互作用 $A\wedge A$ 在一阶变分中不可见，但在二阶变分与 Yang--Mills 作用量中起决定作用。

Yang--Mills 作用量 $S=\int\operatorname{tr}(F\wedge *F)$ 在 $\widetilde\nabla=\nabla+A$ 下展开为
\begin{equation*}
 S(\widetilde\nabla)=S(\nabla)+2\int\operatorname{tr}(D_\nabla A\wedge *F)+2\int\operatorname{tr}(A\wedge A\wedge *F)+\int\operatorname{tr}((D_\nabla A+A\wedge A)\wedge *(D_\nabla A+A\wedge A)).
\end{equation*}
第一阶变分给 Yang--Mills 方程 $D_\nabla *F=0$；二阶变分含 $A\wedge A$ 的二次与三次项，是规范场自能的非线性来源。Abel 情形 $A\wedge A=0$，作用量退化为 Maxwell 二次型，变分方程线性。

典型例子。
\begin{itemize}
\item $U(1)$ 规范理论（Abel）：$A\wedge A=0$（交换子为零），$\widetilde F=F+\dd A$。电磁场强度在规范势平移下只加 $\dd A$，没有自相互作用——这正是 Maxwell 方程线性的根源。
\item $SU(2)$ Yang--Mills：$\widetilde F=F+D_\nabla A+A\wedge A$，$A\wedge A\neq0$ 产生非 Abel 自相互作用。Yang--Mills 方程 $D_\mu F^{\mu\nu}=0$ 展开后含 $A_\mu A_\nu A^\nu$ 三次项，是非线性性的来源。
\item Riemann 几何：取 $\nabla$ 为 Levi--Civita 联络、$A$ 为 $(1,2)$ 张量，$\widetilde R=R+\nabla A+A\circ A$ 给出两个度量相容联络的曲率差。若 $\widetilde\nabla$ 也无挠，$A$ 对下指标对称，二次项控制如 Rauch 比较定理中的曲率比较。
\end{itemize}

这一展开是规范理论全部相互作用结构的微缩版：线性项 $D_\nabla A$ 是"外源激发"，二次项 $A\wedge A$ 是"自相互作用"。在 Yang--Mills 中，规范场既被外源激发又自我激发，正是非 Abel 规范理论区别于 Maxwell 理论的本质。在广义相对论中，Einstein 张量 $G_{\mu\nu}$ 含 Ricci 曲率，而 Ricci 在联络扰动下的二次项对应引力场的自相互作用——这是 Einstein 方程非线性的几何根源。
\end{proof}

\begin{theorem}[Cartan 结构方程]
设 $\theta^a$ 为切丛余标架，$\omega^a{}_b$ 为联络形式。挠率与曲率二形式分别满足 $T^a=\dd\theta^a+\omega^a{}_b\wedge\theta^b$、$\Omega^a{}_b=\dd\omega^a{}_b+\omega^a{}_c\wedge\omega^c{}_b$。
\end{theorem}
\begin{proof}
先从挠率定义出发。对任意 $X,Y$，
\[
\begin{aligned}
T^a(X,Y)
={}&\theta^a(\nabla_XY-\nabla_YX-[X,Y])\\
={}&X[\theta^a(Y)]-Y[\theta^a(X)]-\theta^a([X,Y])\\
&+\omega^a{}_b(X)\theta^b(Y)
-\omega^a{}_b(Y)\theta^b(X).
\end{aligned}
\]
第一行中前三项恰是 $\dd\theta^a(X,Y)$，后二项恰是
$(\omega^a{}_b\wedge\theta^b)(X,Y)$。故
\begin{equation}\label{eq:math-dg-synthesis-cartan-first}
{T^a=\dd\theta^a+\omega^a{}_b\wedge\theta^b.}
\end{equation}

再对基向量 $e_b$ 展开曲率定义：
\[
\begin{aligned}
R(X,Y)e_b
={}&\nabla_X\bigl(\omega^c{}_b(Y)e_c\bigr)
-\nabla_Y\bigl(\omega^c{}_b(X)e_c\bigr)
-\omega^c{}_b([X,Y])e_c\\
={}&\Bigl(
X[\omega^a{}_b(Y)]-Y[\omega^a{}_b(X)]-\omega^a{}_b([X,Y])
\Bigr)e_a\\
&+\Bigl(
\omega^c{}_b(Y)\omega^a{}_c(X)
-\omega^c{}_b(X)\omega^a{}_c(Y)
\Bigr)e_a.
\end{aligned}
\]
把第一括号认作 $\dd\omega^a{}_b(X,Y)$，第二括号认作
$(\omega^a{}_c\wedge\omega^c{}_b)(X,Y)$，得到
\begin{equation}\label{eq:math-dg-synthesis-cartan-second}
{
\Omega^a{}_b
=\dd\omega^a{}_b+\omega^a{}_c\wedge\omega^c{}_b.}
\end{equation}

第一结构方程 $T^a=\dd\theta^a+\omega^a{}_b\wedge\theta^b$ 说挠率是「外微分加联络修正」；第二结构方程 $\Omega^a{}_b=\dd\omega^a{}_b+\omega^a{}_c\wedge\omega^c{}_b$ 说曲率是联络的「外微分加自乘」。二者合起来把挠率和曲率完全用标架 $\{\theta^a\}$ 和联络 $\{\omega^a{}_b\}$ 表达，无需坐标。对 Riemann 流形（挠率 $T=0$），第一方程显式解出 $\omega^a{}_b$，将其表为 $\dd\theta^a$ 的函数，第二方程则给出曲率——这是 Cartan 活动标架法计算曲率的标准流程。

取球面 $\mathbb S^2$ 的活动标架 $\theta^1=R\,\dd\theta$、$\theta^2=R\sin\theta\,\dd\varphi$。由 $T=0$ 解第一方程得 $\omega^1{}_2=-\cos\theta\,\dd\varphi$。代入第二方程：$\Omega^1{}_2=\dd\omega^1{}_2=\sin\theta\,\dd\theta\wedge\dd\varphi=(1/R^2)\,\theta^1\wedge\theta^2$，立即读出截面曲率 $K=1/R^2$。整个计算无需 Christoffel 符号。
\end{proof}

\begin{proposition}[曲率的齐次规范变换]
若 $A^\prime=g^{-1}Ag+g^{-1}\dd g$，则 $F^\prime=g^{-1}Fg$。因此任何伴随不变多项式作用于曲率所得的标量值形式均不依赖局部标架。
\end{proposition}
\begin{proof}
规范变换下联络形式的变换律为
\begin{equation*}
 A'=g^{-1}Ag+g^{-1}\dd g.
\end{equation*}
记 $\eta=g^{-1}\dd g$（纯规范），则 $A'=g^{-1}Ag+\eta$。目标是对 $F'=\dd A'+A'\wedge A'$ 直接计算并证明 $F'=g^{-1}Fg$。

由 $\eta=g^{-1}\dd g$ 是纯规范，Maurer--Cartan 方程给出
\begin{equation*}
 \dd\eta+\eta\wedge\eta=0.
\end{equation*}
这是 $\eta$ 的"平坦性"条件，几何意义是 $g$ 是全局定义的光滑函数，其微分可积。

由 $\dd(g^{-1})=-\eta\,g^{-1}$（对 $g\cdot g^{-1}=I$ 求微分），
\begin{equation*}
 \dd(g^{-1}Ag)
 =g^{-1}(\dd A)g
 -\eta\wedge g^{-1}Ag
 -g^{-1}Ag\wedge\eta.
\end{equation*}
后两项是 $\eta$ 与 $g^{-1}Ag$ 的交叉楔积，来自 $\dd(g^{-1})$ 与 $\dd g$ 上的 Leibniz 规则。

展开 $F'=\dd A'+A'\wedge A'$ 后收集所有项：
\begin{itemize}
\item $g^{-1}(\dd A)g$：来自 $\dd(g^{-1}Ag)$。
\item $g^{-1}(A\wedge A)g$：来自 $(g^{-1}Ag)\wedge(g^{-1}Ag)$。
\item $\dd\eta$：来自 $\dd\eta$。
\item $\eta\wedge\eta$：来自 $\eta\wedge\eta$。
\item 交叉项 $\eta\wedge g^{-1}Ag$ 和 $g^{-1}Ag\wedge\eta$：来自 $\dd(g^{-1}Ag)$ 展开的后两项和 $A'\wedge A'$ 的展开。
\end{itemize}
具体地，$A'\wedge A'=(g^{-1}Ag+\eta)\wedge(g^{-1}Ag+\eta)$ 展开为四项：$(g^{-1}Ag)\wedge(g^{-1}Ag)$、$(g^{-1}Ag)\wedge\eta$、$\eta\wedge(g^{-1}Ag)$、$\eta\wedge\eta$。其中两个交叉项与 $\dd(g^{-1}Ag)$ 中的两个交叉项符号相反，精确配对抵消。

由 Maurer--Cartan 恒等式 $\dd\eta=-\eta\wedge\eta$，以及交叉项精确配对抵消，只剩
\begin{equation*}
 F'=g^{-1}(\dd A+A\wedge A)g
 =g^{-1}Fg.
\end{equation*}
交叉项的相消可显式核对：$\dd(g^{-1}Ag)$ 给 $-\eta\wedge g^{-1}Ag-g^{-1}Ag\wedge\eta$，$A'\wedge A'$ 给 $+g^{-1}Ag\wedge\eta+\eta\wedge g^{-1}Ag$，二者符号相反；$\dd\eta$ 与 $\eta\wedge\eta$ 则由 Maurer--Cartan 恒等式消去。

典型例子。
\begin{itemize}
\item $U(1)$ 规范：$g=\ee^{\ii\chi}$，$F'=F$（Abel 群共轭平凡）。电磁场强度 $F_{\mu\nu}=\partial_\mu A_\nu-\partial_\nu A_\mu$ 在规范变换 $A_\mu\mapsto A_\mu+\partial_\mu\chi$ 下不变，正是 Maxwell 理论规范不变性的几何根源。
\item $SU(N)$ Yang--Mills：$F'=g^{-1}Fg$，规范场强做共轭变换。$\operatorname{tr}(F\wedge F)$ 因此规范不变，是 Chern--Weil 理论构造第二 Chern 类 $c_2(E)$ 的基础。
\item $SO(n)$ 标架丛：$\Omega'=g^{-1}\Omega g$ 给 Riemann 曲率二形式的齐次变换律，与联络一形式的非齐次律对比鲜明。
\item 仿射联络：$\Gamma$ 在坐标变换下非齐次（含二阶导数项），但 Riemann 曲率张量 $R^\rho{}_{\sigma\mu\nu}$ 齐次变换，是真正的张量。
\item 瞬子解：$SU(2)$ BPST 瞬子 $A_\mu=(1-f)\,g^{-1}\partial_\mu g$ 中 $F=g^{-1}F_0g$ 共轭不变，$\operatorname{tr}(F\wedge F)=\operatorname{tr}(F_0\wedge F_0)$ 给瞬子数 $k=1$。
\end{itemize}

这就是"联络 非张量而 曲率 张量"的矩阵形式证明：$A$ 在 规范 变换下有非齐次项 $g^{-1}\dd g$（仿射变换），但 $F$ 只做共轭变换（齐次变换），因此曲率是真正的张量。在 Yang--Mills 理论中，作用量 $\int\operatorname{tr}(F\wedge *F)$ 规范不变正是源于 $F$ 的齐次变换律与 迹 的循环性。在 Chern--Weil 理论中，$\operatorname{tr}(F^k)$ 规范不变且闭，定义了不依赖联络选择的拓扑不变量。这一"非齐次 $\to$ 齐次"的升级是规范理论全部构造性结论（守恒流、规范不变作用量、拓扑荷）的共同根源。
\end{proof}

\begin{theorem}[Bianchi 恒等式]
由同一切丛联络定义的挠率 $T$ 与曲率 $\Omega$ 满足 $DT^a=\Omega^a{}_b\wedge\theta^b$ 和 $D\Omega=0$；第二式也适用于一般向量丛联络。
\end{theorem}
\begin{proof}
由第一结构方程，
\[
T^a=\dd\theta^a+\omega^a{}_b\wedge\theta^b.
\]
用\eqnrefs{eq:math-dg-exterior-covariant-derivative}计算
\[
\begin{aligned}
D T^a
={}&\dd T^a+\omega^a{}_c\wedge T^c\\
={}&\dd\omega^a{}_b\wedge\theta^b
-\omega^a{}_b\wedge\dd\theta^b
+\omega^a{}_c\wedge\dd\theta^c
+\omega^a{}_c\wedge\omega^c{}_b\wedge\theta^b.
\end{aligned}
\]
中间两项在重命名哑指标后相消，于是
\[
D T^a
=\bigl(\dd\omega^a{}_b+\omega^a{}_c\wedge\omega^c{}_b\bigr)\wedge\theta^b
=\Omega^a{}_b\wedge\theta^b.
\]
因此
\begin{equation}\label{eq:math-dg-bianchi-first}
{D T^a=\Omega^a{}_b\wedge\theta^b.}
\end{equation}

对第二结构方程作协变外微分。对一个 $\End(TM)$-值二形式，
\[
(D\Omega)^a{}_b
=\dd\Omega^a{}_b
+\omega^a{}_c\wedge\Omega^c{}_b
-\Omega^a{}_c\wedge\omega^c{}_b.
\]
代入 $\Omega=\dd\omega+\omega\wedge\omega$，并使用 $\dd^2=0$，
\[
\begin{aligned}
D\Omega
={}&\dd(\omega\wedge\omega)
+\omega\wedge(\dd\omega+\omega\wedge\omega)
-(\dd\omega+\omega\wedge\omega)\wedge\omega\\
={}&(\dd\omega\wedge\omega-\omega\wedge\dd\omega)
+\omega\wedge\dd\omega+\omega\wedge\omega\wedge\omega\\
&-\dd\omega\wedge\omega-\omega\wedge\omega\wedge\omega\\
={}&0.
\end{aligned}
\]
故
\begin{equation}\label{eq:math-dg-bianchi-second}
{D\Omega^a{}_b=0.}
\end{equation}
两式也可合写为
\begin{equation}\label{eq:math-dg-synthesis-bianchi}
D T^a=\Omega^a{}_b\wedge\theta^b,
\qquad
D\Omega^a{}_b=0.
\end{equation}

Bianchi 恒等式是曲率的可积条件：$D\Omega=0$ 说明曲率二形式在外协变微分下闭合。这是 Chern--Weil 理论的出发点：$\operatorname{tr}(\Omega^k)$ 为闭形式，其联络无关性由下一定理证明。二维曲面上的 $D\Omega$ 是三形式，因维数自动为零，并不要求 Gauss 曲率为常数。规范理论的 Bianchi 恒等式同样由联络的曲率定义得到；Einstein 张量的 $\nabla_\mu G^{\mu\nu}=0$ 则来自 Levi--Civita 联络第二 Bianchi 恒等式的指标缩并，不能由独立 Yang--Mills 联络的恒等式推出。
\end{proof}

\begin{theorem}[Chern--Weil 形式的联络无关性]
固定一个复向量丛，令 $A_t$ 为光滑联络族、$F_t$ 为曲率。对 $k\ge1$，$\operatorname{tr}(F_t^k)$ 闭合，且 $\partial_t\operatorname{tr}(F_t^k)=k\dd\operatorname{tr}(\dot A_t\wedge F_t^{k-1})$。因此其 de Rham 类不随联络变化。
\end{theorem}
\begin{proof}
对一族联络 $A_t$ 与曲率 $F_t$，把 $\operatorname{tr}(F_t^k)$ 对 $t$ 求导。利用 $\dot F_t=D_t\dot A_t$、Bianchi 恒等式 $D_tF_t=0$ 与 迹 的循环性，把导数写成 恰当形式，从而 de Rham 上同调类与 $t$ 无关。

对
\begin{equation*}
F_t=\dd A_t+A_t\wedge A_t
\end{equation*}
求导（$\dot A_t:=\frac{\dd A_t}{\dd t}$）：
\begin{equation*}
\dot F_t
=\dd\dot A_t+A_t\wedge\dot A_t+
\dot A_t\wedge A_t
=D_t\dot A_t.
\end{equation*}
最后等式是协变外微分的定义 $D_t\dot A_t=\dd\dot A_t+[A_t,\dot A_t]$，其中 $[A_t,\dot A_t]=A_t\wedge\dot A_t+\dot A_t\wedge A_t$（一形式楔积带符号合并）。

由 Leibniz 规则，
\begin{equation*}
\frac{\dd}{\dd t}\operatorname{tr}(F_t^k)
=k\,\operatorname{tr}(\dot F_t\wedge F_t^{k-1})
=k\,\operatorname{tr}(D_t\dot A_t\wedge F_t^{k-1}).
\end{equation*}
因子 $k$ 来自 $F_t^k=F_t\wedge\cdots\wedge F_t$ 中 $\dot F_t$ 可落在 $k$ 个位置之一，迹 的循环性使这 $k$ 项相等。

利用协变 Leibniz 规则和 Bianchi 恒等式，可把 $D_t$ 从 $\dot A_t$ 移到 $F_t^{k-1}$：
\begin{equation*}
D_t(\dot A_t\wedge F_t^{k-1})
=D_t\dot A_t\wedge F_t^{k-1}
-\dot A_t\wedge D_t(F_t^{k-1}),
\end{equation*}
以及第二 Bianchi 恒等式 $D_tF_t=0$（从而 $D_t(F_t^{k-1})=0$），得
\begin{equation*}
D_t(\dot A_t\wedge F_t^{k-1})
=D_t\dot A_t\wedge F_t^{k-1}.
\end{equation*}
这一步是关键：Bianchi 恒等式让 $D_t$ 在乘积上"穿过"所有 $F_t$ 因子，仿佛它们是"协变常量"。

迹把 $D_t$ 退化为 $\dd$。对任意次数为 $p$ 的矩阵值形式 $\beta$，分次循环性给出 $\operatorname{tr}(A_t\wedge\beta)=(-1)^p\operatorname{tr}(\beta\wedge A_t)$，故分次交换子 $A_t\wedge\beta-(-1)^p\beta\wedge A_t$ 的迹为零。因此
\begin{equation*}
\operatorname{tr}(D_t(\dot A_t\wedge F_t^{k-1}))
=\dd\operatorname{tr}(\dot A_t\wedge F_t^{k-1}),
\end{equation*}
得到\eqnrefs{eq:chern-weil-transgression}。

积分 $t\in[0,1]$ 后，
\begin{equation*}
\operatorname{tr}(F_1^k)-\operatorname{tr}(F_0^k)
=k\,\dd\int_0^1\operatorname{tr}(\dot A_t\wedge F_t^{k-1})\,\dd t,
\end{equation*}
两组 联络 产生的 特征 形式 相差 恰当形式，因此代表同一 de Rham 类 $[\operatorname{tr}(F^k)]\in H_{\mathrm{dR}}^{2k}(M)$。

典型例子。
\begin{itemize}
\item $U(1)$ 线丛：$k=1$，$\operatorname{tr}(F)=F$ 是闭二形式，$[\ii F/(2\pi)]\in H^2_{\mathrm{dR}}(M)$ 是第一 Chern 类 $c_1(E)$。两个不同联络的 $F$ 相差 恰当形式，但积分 $\frac{\ii}{2\pi}\int_M F$ 给出 Chern 数（如磁通量子化）。
\item $SU(N)$ 丛：$k=2$，$\operatorname{tr}(F\wedge F)$ 是闭四形式，$[\operatorname{tr}(F\wedge F)/8\pi^2]$ 是第二 Chern 类 $c_2(E)$。Yang--Mills 瞬子数 $k=\int\operatorname{tr}(F\wedge F)/8\pi^2$ 是拓扑不变量，不依赖规范选择。
\end{itemize}

Chern--Weil 理论给出从 $\mathfrak g$ 上的不变多项式到 $H_{\mathrm{dR}}^*(M)$ 的同态：每个 $G$-不变多项式 $P$ 给出闭形式 $P(F)$，其 de Rham 类只依赖丛 $E$ 的拓扑。结合 Atiyah--Singer 指标定理，这些特征类积分给出椭圆算子的解析指标——如 Dirac 算子的 $\hat A$ 类、符号丛的 Chern 字符。这是"曲率局部数据决定全局拓扑"的精确机制。
\end{proof}

\begin{proposition}[路径序指数的收敛构造]
设 $A_t$ 为有限闭区间上的连续矩阵函数。初值问题 $\dot V=-A_tV$ 的解为绝对一致收敛的 Dyson 级数。若所有 $A_t$ 两两对易，则解退化为 $V(t)=\exp(-\int_0^tA_s\,\dd s)V(0)$。
\end{proposition}
\begin{proof}
把\eqnrefs{eq:parallel-matrix-ode} 即 $\dot V(t)=-A_tV(t)$、$V(0)$ 给定，从 $0$ 到 $t$ 积分得积分方程
\begin{equation*}
V(t)-V(0)=-\int_0^tA_{t_1}V(t_1)\,\dd t_1,
\end{equation*}
即
\begin{equation*}
V(t)=V(0)-\int_0^tA_{t_1}V(t_1)\,\dd t_1.
\end{equation*}
把右端 $V(t_1)$ 用同一积分方程代入，$V(t_1)=V(0)-\int_0^{t_1}A_{t_2}V(t_2)\,\dd t_2$，得
\begin{align*}
V(t)
&=V(0)-\int_0^tA_{t_1}
 \Bigl[V(0)-\int_0^{t_1}A_{t_2}V(t_2)\,\dd t_2\Bigr]\dd t_1\\
&=V(0)-\int_0^tA_{t_1}\dd t_1\,V(0)
 +\int_0^t\int_0^{t_1}A_{t_1}A_{t_2}V(t_2)\,\dd t_2\dd t_1\\
&=\Bigl[I-\int_0^tA_{t_1}\dd t_1\Bigr]V(0)
 +\int_0^t\int_0^{t_1}A_{t_1}A_{t_2}V(t_2)\,\dd t_2\dd t_1.
\end{align*}
对余项中的 $V(t_2)$ 再代入一次，$V(t_2)=V(0)-\int_0^{t_2}A_{t_3}V(t_3)\,\dd t_3$，
\begin{align*}
\int_0^t\int_0^{t_1}A_{t_1}A_{t_2}V(t_2)\,\dd t_2\dd t_1
&=\int_0^t\int_0^{t_1}A_{t_1}A_{t_2}\dd t_2\dd t_1\,V(0)\\
&\quad-\int_0^t\int_0^{t_1}\int_0^{t_2}
 A_{t_1}A_{t_2}A_{t_3}V(t_3)\,\dd t_3\dd t_2\dd t_1.
\end{align*}
合并得第三次迭代
\begin{align*}
V(t)
&=\Bigl[I-\int_0^tA_{t_1}\dd t_1
 +\int_0^t\int_0^{t_1}A_{t_1}A_{t_2}\dd t_2\dd t_1\Bigr]V(0)\\
&\quad-\int_0^t\int_0^{t_1}\int_0^{t_2}
 A_{t_1}A_{t_2}A_{t_3}V(t_3)\,\dd t_3\dd t_2\dd t_1.
\end{align*}
由归纳法，截断至 $N$ 阶的有限展开为
\begin{equation*}
V(t)=\sum_{j=0}^{N}(-1)^j
\int_{0<t_j<\cdots<t_1<t}
A_{t_1}\cdots A_{t_j}\,\dd t_j\cdots\dd t_1\,V(0)+R_{N+1}(t),
\end{equation*}
其中 $j=0$ 项理解为 $I\,V(0)$，余项为
\[
R_{N+1}(t)=(-1)^{N+1}\int_{0<t_{N+1}<\cdots<t_1<t}
A_{t_1}\cdots A_{t_{N+1}}V(t_{N+1})\,\dd t_{N+1}\cdots\dd t_1.
\]
在有限区间 $0\le t\le T$ 上令 $M=\sup_t\|A_t\|$、$B=\sup_t\|V(t)\|<\infty$，则单纯形体积给出
\[
\|R_{N+1}(t)\|\le B\frac{M^{N+1}T^{N+1}}{(N+1)!}\longrightarrow0,
\qquad
\|\text{第 }j\text{ 阶项}\|\le\|V(0)\|\frac{M^jT^j}{j!}.
\]
这分别证明余项一致趋零及级数绝对一致收敛，从而产生 Dyson 级数
\begin{equation*}
V(t)=\Bigl[I+\sum_{k=1}^\infty(-1)^k
\int_{0<t_k<\cdots<t_1<t}
A_{t_1}\cdots A_{t_k}\,\dd t_k\cdots\dd t_1\Bigr]V(0).
\end{equation*}
若所有 $A_t$ 对易，积分区域的 $k!$ 个排列给相同 integrand，于是有序单纯形积分等于立方体积分的 $1/k!$，级数便退化为普通指数。这解释了什么时候可以安全地把 $\mathcal P$ 去掉。
\end{proof}
